%=============================================================================
% MODULE 06: RESULTS
%=============================================================================
% Frontiers in Public Health — Objective results, no interpretation
%   6.1 Model performance
%   6.2 Fairness analysis
%   6.3 Intersectional analysis
%   6.4 Cross-validation stability
%=============================================================================

\section{Results}

\subsection{Model Performance}

Table~\ref{tab:performance} presents the discriminative performance of 
all four models on the held-out test set.

\begin{table}[H]
\centering
\caption{Performance Comparison of Machine Learning Models}
\begin{tabular}{lcccc}
\toprule
\textbf{Model} & \textbf{CV F1 (mean±sd)} & \textbf{Test F1} & 
\textbf{AUC-ROC} & \textbf{Accuracy} \\
\midrule
Logistic Regression & 0.7929±0.0313 & 0.8272 & 0.8877 & 0.8052 \\
Random Forest & 0.7679±0.0243 & 0.7875 & 0.8731 & 0.7727 \\
Gradient Boosting & 0.7769±0.0221 & 0.7805 & 0.8543 & 0.7792 \\
\textbf{SVM} & \textbf{0.7859±0.0229} & \textbf{0.8313} & \textbf{0.8794} & \textbf{0.8117} \\
\bottomrule
\end{tabular}
\label{tab:performance}
\end{table}

The SVM model achieved the highest F1-Score (0.8313) and competitive 
AUC-ROC (0.8794), with stable cross-validation performance (CV: 
0.7859$\pm$0.0229). Logistic Regression demonstrated the best AUC-ROC 
(0.8877) but lower F1-Score (0.8272), suggesting a trade-off between 
discrimination and class-specific performance. Random Forest and Gradient 
Boosting showed comparable but slightly lower performance across all metrics.

\subsection{Fairness Analysis}

\subsubsection{Ethnicity-Based Fairness}

Table~\ref{tab:ethnicity_fairness} presents fairness metrics stratified by 
ethnicity for all models.

\begin{table}[H]
\centering
\caption{Fairness Metrics by Ethnicity}
\begin{tabular}{lcccccc}
\toprule
\textbf{Model} & \textbf{EOD} & \textbf{DPD} & 
\textbf{TPR$_{\text{maj}}$} & \textbf{TPR$_{\text{min}}$} & 
\textbf{FPR$_{\text{maj}}$} & \textbf{FPR$_{\text{min}}$} \\
\midrule
Log.\ Regression & 0.0687 & 0.0479 & 0.833 & 0.815 & 0.182 & 0.375 \\
Random Forest & 0.0306 & 0.0324 & 0.833 & 0.778 & 0.200 & 0.417 \\
Grad.\ Boosting & 0.0636 & 0.0201 & 0.812 & 0.704 & 0.164 & 0.333 \\
\textbf{SVM} & \textbf{0.0492} & \textbf{0.0148} & 0.812 & 0.741 & 0.182 & 0.375 \\
\bottomrule
\end{tabular}
\label{tab:ethnicity_fairness}
\end{table}

All models exhibited higher false positive rates for minority groups 
(FPR: 0.333--0.417) compared to majority groups (FPR: 0.164--0.200). 
The Random Forest model showed the lowest EOD (0.0306) for ethnicity, 
while the SVM model achieved the lowest DPD (0.0148), indicating 
different models exhibit different fairness profiles depending on the 
metric used.

\subsubsection{Gender-Based Fairness}

Table~\ref{tab:gender_fairness} presents fairness metrics stratified by gender.

\begin{table}[H]
\centering
\caption{Fairness Metrics by Gender}
\begin{tabular}{lcccccc}
\toprule
\textbf{Model} & \textbf{EOD} & \textbf{DPD} & 
\textbf{TPR$_{\text{M}}$} & \textbf{TPR$_{\text{F}}$} & 
\textbf{FPR$_{\text{M}}$} & \textbf{FPR$_{\text{F}}$} \\
\midrule
Log.\ Regression & 0.0838 & 0.1010 & 0.833 & 0.815 & 0.182 & 0.375 \\
Random Forest & 0.1143 & 0.1010 & 0.833 & 0.778 & 0.200 & 0.417 \\
Grad.\ Boosting & 0.0787 & 0.0843 & 0.812 & 0.704 & 0.164 & 0.333 \\
\textbf{SVM} & \textbf{0.0566} & \textbf{0.0672} & 0.812 & 0.741 & 0.182 & 0.375 \\
\bottomrule
\end{tabular}
\label{tab:gender_fairness}
\end{table}

Gender-based analysis revealed similar patterns, with female patients 
experiencing higher false positive rates across all models. The SVM model 
achieved the lowest EOD (0.0566) and DPD (0.0672) for gender, indicating 
the most equitable performance across this dimension.

\subsection{Intersectional Analysis}

Table~\ref{tab:intersectional} presents performance metrics at the 
intersection of gender and ethnicity.

\begin{table}[H]
\centering
\caption{Intersectional Analysis: Performance by Gender $\times$ Ethnicity}
\begin{tabular}{lcccc}
\toprule
\textbf{Group} & \textbf{Accuracy} & \textbf{F1-Score} & 
\textbf{TPR} & \textbf{$n$} \\
\midrule
Majority--Female & 0.8000 & 0.8235 & 0.833 & 65 \\
Majority--Male & 0.8125 & 0.8276 & 0.833 & 38 \\
Minority--Female & 0.7826 & 0.7500 & 0.714 & 33 \\
Minority--Male & 0.8000 & 0.7273 & 0.750 & 18 \\
\bottomrule
\end{tabular}
\label{tab:intersectional}
\end{table}

The intersectional analysis revealed that minority--female patients 
experienced the lowest TPR (0.714) and F1-Score (0.750), indicating 
amplified bias at the intersection of multiple marginalized identities. 
This group also showed the largest performance gap relative to the 
best-performing group (majority--male), with a TPR difference of 0.119 
and an F1-Score difference of 0.078.

\subsection{Cross-Validation Stability}

Table~\ref{tab:cross_validation} summarizes cross-validation results 
across all models.

\begin{table}[H]
\centering
\caption{Cross-Validation Stability Assessment}
\begin{tabular}{lccc}
\toprule
\textbf{Model} & \textbf{Mean F1} & \textbf{SD F1} & \textbf{CV Folds} \\
\midrule
Logistic Regression & 0.7929 & 0.0313 & 5 \\
Random Forest & 0.7679 & 0.0243 & 5 \\
Gradient Boosting & 0.7769 & 0.0221 & 5 \\
SVM & 0.7859 & 0.0229 & 5 \\
\bottomrule
\end{tabular}
\label{tab:cross_validation}
\end{table}

All models demonstrated stable performance across cross-validation folds, 
with standard deviations below 0.04. The Gradient Boosting model showed 
the highest stability (SD: 0.0221), followed closely by SVM (SD: 0.0229). 
These results indicate that the observed performance and fairness metrics 
are reliable and not attributable to random variation in data splitting.
